% RESEARCH_DRAFT_SHA256: acc87b1482a028def4bf7fc427140396b5f7d308e2f17f01d81b356b73432925
% Layout and source-reference wrappers expanded for the web renderer.
\documentclass[11pt]{article}
\usepackage[margin=25mm]{geometry}
\usepackage{amsthm}
\usepackage{fontspec,unicode-math}
\setmainfont{Latin Modern Roman}
\setmathfont{Latin Modern Math}
\usepackage{amsmath,amssymb,mathtools}
\usepackage{xurl,hyperref}
\hypersetup{colorlinks=true,urlcolor=blue,linkcolor=blue}
\setcounter{secnumdepth}{0}
\setlength{\parindent}{0pt}
\setlength{\parskip}{5pt}
\setlength{\emergencystretch}{3em}
\newcommand{\sref}[2]{\hyperlink{source:#1:#2}{[S#2]}}
\newcommand{\cataloguescope}[1]{\paragraph{Source question.}#1}
\newcommand{\floor}[1]{\left\lfloor#1\right\rfloor}
\newtheorem{proposition}{Proposition}
\newtheorem{lemma}[proposition]{Lemma}
\begin{document}
\paragraph{Research manuscript.} OpenAI \texttt{gpt-6-astra}, reasoning \texttt{ultra}; 11 October 2026.
\paragraph{Preserved source body.} The complete catalogue mathematical text and sources follow with layout and source links expanded, including source qualifications and its historical preparation notes. New research begins at the section headed Research attempt.

\section{Integral Primal Certificates for Submodular Minimization}\label{statement}

\subsection{Definitions and mathematical statement}

\paragraph{Shared notation.}Standard number systems are denoted by $\mathbb N,\mathbb Z,\mathbb Q,\mathbb R,\mathbb C$; index ranges, computational models and specialized notation follow the cited problem source and the definitions below.


For normalized $f(\varnothing)=0$, its base polyhedron is $B(f)=\{y\in\mathbb R^V:y(S)\leq f(S)\ (S\subseteq V),\ y(V)=f(V)\}$ with $y(S)=\sum_{e\in S}y_e$. The associated primal objective is $\max_{y\in B(f)}\sum_e\min(y_e,0)$. Greedy vertices take successive value differences along permutations. A convex combination is $y=\sum_i\lambda_i v_i$ with $\lambda_i\geq0$ and $\sum_i\lambda_i=1$.

A set function $f:2^V\to\mathbb R$ is submodular if $f(A)+f(B)\geq f(A\cup B)+f(A\cap B)$ for all $A,B\subseteq V$. It is monotone if inclusion never decreases its value and normalized if $f(\varnothing)=0$.

A polyhedron is an intersection of finitely many closed affine half-spaces; a polytope is a bounded polyhedron. Its faces are its intersections with supporting hyperplanes.
\paragraph{Statement recorded in the supplied source.}
For submodular function minimization, find an algorithm that directly produces an integral primal optimal point in the base polyhedron together with its representation as a convex combination of greedy vertices. \href{https://doi.org/10.4171/owr/2005/50}{[S2]}
\paragraph{Scope qualification.}
The report uses real-valued $f$ but its requested integral optimum requires appropriate integrality assumptions. These are not silently imposed on the original question. The normalization above states the convention needed by the displayed base representation.
\paragraph{Source status.}
The supplied catalogue labels this record open. Its source-status metadata was checked against the supplied snapshot on 2026-10-11. This import does not independently establish current open status or mathematical correctness. The catalogue literature assessment is dated 2026-08-21.
\subsection{Short English statement}
Produce the requested integral primal optimum and an explicit decomposition into greedy vertices for submodular minimization.
\subsection{Sources}
\begin{itemize}
\item[S1] Ulam AI and the UnsolvedMath Contributors, UnsolvedMath: A Curated Collection of Open Mathematics Problems, catalogue entry. Catalogue curation and metadata: CC BY 4.0. The supplied source question is retained; definitions and the English summary are editorial additions.
\url{https://huggingface.co/datasets/ulamai/UnsolvedMath}.
\item[S2] Combinatorial Optimization (2006), Oberwolfach Reports, DOI 10.4171/owr/2005/50, pp. 2833--2932
\url{https://doi.org/10.4171/owr/2005/50}.
\item[S3] Original Oberwolfach Report 50/2005, pp. 2884--2885
\url{https://ems.press/content/serial-article-files/46024}.
\textit{Underlying Oberwolfach Reports material retains its source-specific licence, including CC BY-SA 4.0 where applicable; see THIRD\_PARTY\_NOTICES.md.}
\end{itemize}
\paragraph{Preparation and review.}
Original question attributed to the cited source and UnsolvedMath contributors. Codex (GPT-6; exact serving revision not exposed) assisted with definitions, English summaries, source inspection and formatting. No new human mathematical authorship or independent proof verification is asserted. Catalogue identity checked against the complete supplied registry and existing repository statements; known flagged repetitions and obvious identical conjectures excluded. Full mathematical-equivalence/literature review has not been performed. Statement preparation and final review followed contribution prompts C and D; automated checks verify structure and identity, not open status.
\label{end:statement}

\section{Research attempt}
\subsection{Literature, source context and scope}
The full original catalogue body and source qualifications are preserved above. The relevant primary workshop passages were read in the cached publisher report. Historical claims and subsequent attributed results are distinguished from the mathematical work below. A restricted example, standard lemma or known theorem is not counted as new progress on the entire original target. Missing assumptions and repeated or rhetorical extracts are not silently turned into different solved problems.

\subsection{Testing integrality assumptions and greedy-vertex search}
The source states a submodular minimization duality involving the base polyhedron, then asks for a direct construction of an integral dual optimum and its greedy decomposition. For a normalized submodular function $f$ on $E$, use
\[
B(f)=\{y\in\mathbb R^E:y(S)\le f(S)\ (S\subseteq E),\ y(E)=f(E)\},
\qquad g(y)=\sum_{e\in E}\min(y_e,0).
\]
If $f$ is allowed to be arbitrary real valued, an integral base point need not exist: on a singleton set take $f(\varnothing)=0,f(\{1\})=1/2$. Submodularity is immediate, and $B(f)=\{1/2\}$. This only identifies a missing hypothesis in the literal extraction; it is not a counterexample to an intended integer-valued theorem.
\begin{proposition}
Even for integer-valued submodular $f$, maximizing $g$ by inspecting only greedy vertices can fail.
\end{proposition}
\begin{proof}
Take $E=\{1,2\}$, $f(\varnothing)=f(E)=0$ and $f(\{1\})=f(\{2\})=1$. The only nonnested submodularity inequality is $1+1\ge0+0$, and nested pairs give equality, so $f$ is submodular. Its base polyhedron is precisely $y=(t,-t)$ with $-1\le t\le1$: the singleton upper bounds and zero total give these inequalities. The two greedy orders have increments $(1,-1)$ and $(-1,1)$, respectively. On the segment,
\[
g(t,-t)=\min(t,0)+\min(-t,0)=-|t|.
\]
The unique maximizer is the integral point $(0,0)$ with value zero, whereas both greedy vertices have value $-1$. The optimum has the explicit convex decomposition $\frac12(1,-1)+\frac12(-1,1)$.
\end{proof}
\subsection{Verification and remaining gaps}
All submodularity inequalities and both greedy orders on the two-element example were checked exactly. The failed inference is that a maximizer must be a vertex because the feasible region is a polyhedron: here the objective is concave piecewise linear rather than linear, and its maximum lies in the interior of an edge. Convex combinations of greedy bases are therefore essential even in this small case. No direct polynomial algorithm or general integral-optimum decomposition procedure is supplied. The source already gives duality and relevant prior algorithms; neither that literature nor the example is claimed as a new solution.

\section{References and source attribution}
Primary workshop report: \url{https://doi.org/10.4171/owr/2005/50}. Relevant passages read: 2005/50 pp. 2883--2885, submodular minimization duality, greedy bases and the direct-construction question.

All retrievals were checked on 11 October 2026. Original question attribution remains in the preserved Sources section. External theorems are attributed and not claimed as new discoveries or fully verified merely by citation.


\paragraph{Integration review.} Codex (GPT-6; exact serving revision not exposed) reviewed the written arguments and expanded layout and source-reference wrappers. Repository metadata and the production web renderer were checked using Python 3.13 and JavaScript with MathJax 3.2.2. This model-assisted review is not independent human proof verification; no PDF compilation or formal proof checking is asserted.

\section{Completion Estimate}
\noindent\textbf{COMPLETION: 0\%}
No substantiated new progress toward the entire original target is claimed. The worked results and failed route are delimited above; a known result, illustrative model, or auxiliary obstruction is not counted as original-target completion.

\section{Final status and attribution}
\textbf{Attempt status: unresolved.} This describes this research attempt, not a new ruling on the source catalogue's status or on external literature claims.
The mathematics and this self-audit were generated by OpenAI \texttt{gpt-6-astra}, reasoning \texttt{ultra}, on 11 October 2026. Codex (GPT-6; exact serving revision not exposed) orchestrated the batch and contributed formatting/repository review. No human mathematical authorship, independent human proof checking, or formal verification is claimed. The historical preparation paragraphs in the reproduced catalogue body concern that body only.

Contribution prompt A was applied to the complete definition and supplied source record, followed by the adversarial checks of prompt B. Local reproduction files are \texttt{batch016/checks.py}, \texttt{computations.json}, and the per-ID audit JSON. They were not public links at preparation time. Numerical checks are finite and do not prove unrestricted results. No LaTeX engine was found locally; compilation was not performed by this research agent.

\end{document}
